\section{Week 5 --- Capacitance, Electrostatic Energy, and Multipole Expansion}

\subsection{Electrostatic Systems of Conductors and Capacitance}

\begin{definition}{Electrostatic Capacitance}{capacitance}
Consider two isolated conductors embedded in a linear, homogeneous, and isotropic dielectric medium (or vacuum) carrying net charges \(+Q\) and \(-Q\), respectively. Since each conducting body forms an equipotential volume in electrostatics (\(\phi = \text{constant}\)), the electric potential difference between the positive conductor (\(\phi_+\)) and the negative conductor (\(\phi_-\)) is unambiguously defined along any path from \((-)\) to \((+)\):
\begin{equation}
    V \equiv \phi_+ - \phi_- = -\int_{(-)}^{(+)} \bm{E} \cdot \dd{\bm{s}}.
\end{equation}
The \textbf{capacitance} \(C\) of the two-conductor configuration is the scalar ratio:
\begin{equation}
    C \equiv \frac{Q}{V}.
\end{equation}
Capacitance is an intrinsically positive quantity (\(C > 0\)), whose SI unit is the Farad (\(\text{F} \equiv \text{C}\cdot\text{V}^{-1}\)). In practical applications, capacitances typically scale within the ranges of picofarads (\(\text{pF} = 10^{-12}\,\text{F}\)) to microfarads (\(\mu\text{F} = 10^{-6}\,\text{F}\)).
\end{definition}

\begin{proposition}{Geometric Invariance and Linearity of Capacitance}{prop_cap_geom}
The capacitance \(C\) depends strictly on the geometry of the system---specifically the spatial shapes, dimensions, and mutual separations of the conductors---and is completely independent of the applied potential difference \(V\) or the net stored charge \(Q\).
\end{proposition}

\begin{proofbox}[Justification via Linearity and the Second Uniqueness Theorem]
Coulomb's law dictates that the electrostatic field generated by a charge density distribution \(\rho(\gr')\) is linear with respect to the source charge:
\begin{equation}
    \bm{E}(\gr) = \frac{1}{4\pi\epsilon_0} \int_V \frac{\rho(\gr')}{|\gr - \gr'|^2}\,\hat{\bm{r}}^*\,\dd{V'}.
\end{equation}
Scaling the net charge on each conductor by an arbitrary real factor \(\alpha\) (\(Q \to \alpha Q\)) scales the surface charge distribution \(\sigma(\gr')\) proportionally (\(\sigma \to \alpha \sigma\)). By the \textbf{Second Uniqueness Theorem of Electrostatics}, the electric field \(\bm{E}(\gr)\) is uniquely determined everywhere in the exterior domain once the total charge on each isolated conductor is specified. Consequently, \(\bm{E} \to \alpha \bm{E}\). 

Because the potential difference is the line integral of \(\bm{E}\), it satisfies:
\begin{equation}
    V \to -\int_{(-)}^{(+)} (\alpha \bm{E}) \cdot \dd{\bm{s}} = \alpha V.
\end{equation}
The ratio \(C = Q/V = (\alpha Q)/(\alpha V)\) is therefore invariant under variations of charge or voltage, demonstrating that \(C\) is a purely geometric property of the configuration.
\end{proofbox}

\begin{definition}{Capacitance of an Isolated Conductor}{isolated_capacitance}
For a single isolated conducting body carrying a net charge \(Q\), the second conductor is conventionally defined as an infinitely distant, grounded spherical shell (\(\phi_\infty = 0\)). The capacitance of the isolated body is:
\begin{equation}
    C \equiv \frac{Q}{\phi - \phi_\infty} = \frac{Q}{\phi}.
\end{equation}
\end{definition}

\subsubsection{Canonical Geometries and Exact Capacitance Calculations}

\begin{example}[Isolated Conducting Sphere]
Consider a solid conducting sphere of radius \(R\) carrying a total charge \(Q\) in vacuum. Construct a concentric spherical Gaussian surface of radius \(r \ge R\). By spherical symmetry, Gauss's law yields:
\begin{equation}
    \oint_{\partial V} \bm{E} \cdot \dd{\bm{a}} = E_r (4\pi r^2) = \frac{Q}{\epsilon_0} \implies E_r(r) = \frac{Q}{4\pi\epsilon_0 r^2}.
\end{equation}
The electric potential at the surface relative to reference at infinity is:
\begin{equation}
    V = \phi(R) - \phi(\infty) = -\int_\infty^R E_r\,\dd{r} = \frac{Q}{4\pi\epsilon_0 R}.
\end{equation}
The self-capacitance of the isolated sphere is therefore:
\begin{equation}
    C_{\text{sphere}} = \frac{Q}{V} = 4\pi\epsilon_0 R.
\end{equation}
\end{example}

\begin{example}[Parallel-Plate Capacitor]
Consider two planar parallel conducting plates of identical surface area \(A\), separated by a gap \(d\), carrying charges \(+Q\) and \(-Q\). Assuming the plates are wide compared to their separation (\(A \gg d^2\)), fringing edge fields are negligible.

\begin{center}
\begin{tikzpicture}[scale=0.9, >=Stealth]
    \draw[thick, fill=blue!10] (0, 2) rectangle (6, 2.3);
    \draw[thick, fill=blue!10] (0, 0) rectangle (6, -0.3);
    \node at (3, 2.15) {\small \(+Q \quad (\sigma = +Q/A)\)};
    \node at (3, -0.15) {\small \(-Q \quad (\sigma = -Q/A)\)};
    \foreach \x in {0.8, 1.8, 2.8, 3.8, 4.8, 5.8} {
        \draw[->, thick, red!80!black] (\x, 2.0) -- (\x, 0.0);
    }
    \node[red!80!black] at (6.4, 1.0) {\(\bm{E}\)};
    \draw[<->, thick] (-0.5, 0) -- (-0.5, 2) node[midway, left] {\(d\)};
    \draw[dashed, ForestGreen, thick] (4.0, 1.5) rectangle (5.2, 2.15);
    \node[ForestGreen, right] at (5.2, 1.8) {\footnotesize Gaussian pillbox \(A'\)};
\end{tikzpicture}
\end{center}

Applying Gauss's law to a pillbox of cross-sectional area \(A'\) enclosing the inner surface of the upper plate gives:
\begin{equation}
    \oint \bm{E} \cdot \dd{\bm{a}} = E A' = \frac{Q_{\text{enc}}}{\epsilon_0} = \frac{\sigma A'}{\epsilon_0} \implies E = \frac{\sigma}{\epsilon_0} = \frac{Q}{\epsilon_0 A}.
\end{equation}
The potential difference across the gap is:
\begin{equation}
    V = -\int_{0}^{d} (-E)\,\dd{z} = E d = \frac{Q d}{\epsilon_0 A}.
\end{equation}
Thus, the capacitance is:
\begin{equation}
    C = \frac{Q}{V} = \frac{\epsilon_0 A}{d}.
\end{equation}
\end{example}

\begin{example}[Cylindrical Coaxial Capacitor]
Consider a solid inner cylindrical conductor of radius \(a\) surrounded coaxially by a cylindrical conducting shell of inner radius \(b\), both of length \(L \gg (b - a)\). The conductors carry total charges \(+Q\) and \(-Q\), corresponding to linear charge density \(\lambda = Q/L\).

\begin{enumerate}
    \item \textbf{Electric Field via Gauss's Law}: For an evaluation cylinder of radius \(r \in [a, b]\) and length \(l\):
    \begin{equation}
        \oint \bm{E} \cdot \dd{\bm{a}} = E_r(r)(2\pi r l) = \frac{\lambda l}{\epsilon_0} \implies \bm{E}(r) = \frac{\lambda}{2\pi\epsilon_0 r}\,\hat{\bm{r}}.
    \end{equation}
    By charge neutrality and symmetry, \(\bm{E} = \bm{0}\) for \(r < a\) and \(r > b\).
    \item \textbf{Potential Difference}:
    \begin{equation}
        V = \phi(a) - \phi(b) = -\int_b^a E_r(r)\,\dd{r} = \frac{\lambda}{2\pi\epsilon_0} \int_a^b \frac{\dd{r}}{r} = \frac{\lambda}{2\pi\epsilon_0}\ln\left(\frac{b}{a}\right).
    \end{equation}
    \item \textbf{Capacitance}:
    \begin{equation}
        C = \frac{Q}{V} = \frac{\lambda L}{\frac{\lambda}{2\pi\epsilon_0}\ln(b/a)} = \frac{2\pi\epsilon_0 L}{\ln(b/a)}.
    \end{equation}
\end{enumerate}
\end{example}

\begin{example}[Concentric Spherical Capacitor]
Consider two concentric spherical conducting shells of radii \(a\) and \(b\) (\(a < b\)), carrying charges \(+Q\) and \(-Q\).
\begin{enumerate}
    \item \textbf{Radial Electric Field}: Gauss's law yields \(\bm{E}(r) = \frac{Q}{4\pi\epsilon_0 r^2}\hat{\bm{r}}\) for \(a < r < b\), and \(\bm{E} = \bm{0}\) elsewhere.
    \item \textbf{Potential Difference}:
    \begin{equation}
        V = \phi(a) - \phi(b) = -\int_b^a \frac{Q}{4\pi\epsilon_0 r^2}\,\dd{r} = \frac{Q}{4\pi\epsilon_0}\left(\frac{1}{a} - \frac{1}{b}\right) = \frac{Q}{4\pi\epsilon_0}\frac{b - a}{a b}.
    \end{equation}
    \item \textbf{Capacitance}:
    \begin{equation}
        C = \frac{Q}{V} = 4\pi\epsilon_0 \frac{ab}{b - a}.
    \end{equation}
    Taking the asymptotic limit as the outer shell recedes to infinity (\(b \to \infty\)):
    \begin{equation}
        \lim_{b \to \infty} C = \lim_{b \to \infty} 4\pi\epsilon_0 \frac{a}{1 - a/b} = 4\pi\epsilon_0 a,
    \end{equation}
    which recovers identically the self-capacitance of an isolated sphere of radius \(a\).
\end{enumerate}
\end{example}

\subsection{Electrostatic Energy and Field Energy Density}

\subsubsection{Work of Charging a Capacitor}

\begin{proposition}{Stored Electrostatic Potential Energy}{prop_stored_energy}
The total work \(W\) required by an external agent to transfer charge against the prevailing electrostatic potential gradient from one plate to the other until a final charge state \((\pm Q)\) is attained is stored as electrostatic potential energy \(U\):
\begin{equation}
    U = \frac{1}{2}\frac{Q^2}{C} = \frac{1}{2}QV = \frac{1}{2}CV^2.
\end{equation}
\end{proposition}

\begin{proofbox}[Quasi-static Charging Process]
Let \(q(t)\) denote the instantaneous charge separated between the plates at some point during the charging process. The instantaneous potential difference is \(v(q) = q/C\). The incremental external mechanical work \(\dd{W}\) required to transfer an additional infinitesimal charge element \(\dd{q}\) from the negative plate to the positive plate is:
\begin{equation}
    \dd{W} = v(q)\,\dd{q} = \frac{q}{C}\,\dd{q}.
\end{equation}
Integrating from an initial uncharged state \(q = 0\) to the final charge \(q = Q\):
\begin{equation}
    W = \int_0^Q \frac{q}{C}\,\dd{q} = \frac{1}{2}\frac{Q^2}{C}.
\end{equation}
Using the relation \(Q = CV\), the alternative expressions \(U = \frac{1}{2}QV = \frac{1}{2}CV^2\) follow directly.
\end{proofbox}

\subsubsection{Energy Density of the Electrostatic Field}

Expressing the capacitance and voltage of a parallel-plate capacitor in terms of the uniform intermediate field \(\bm{E}\) (\(C = \epsilon_0 A / d\) and \(V = E d\)):
\begin{equation}
    U = \frac{1}{2} C V^2 = \frac{1}{2} \left(\frac{\epsilon_0 A}{d}\right) (E d)^2 = \frac{1}{2}\epsilon_0 E^2 (A d).
\end{equation}
Since \(A d = \text{Volume}\) occupied by the field between the plates, we localize the electrostatic potential energy within the field itself.

\begin{definition}{Electrostatic Energy Density}{energy_density}
The electrostatic energy stored per unit volume of space occupied by an electric field \(\bm{E}\) in vacuum is defined by:
\begin{equation}
    u_E \equiv \frac{\dd{U}}{\dd{V}} = \frac{1}{2}\epsilon_0 |\bm{E}|^2.
\end{equation}
The total electrostatic energy contained within an arbitrary spatial volume \(\mathcal{V}\) is:
\begin{equation}
    U = \frac{\epsilon_0}{2} \int_{\mathcal{V}} |\bm{E}(\gr)|^2\,\dd{V}.
\end{equation}
\end{definition}

\subsubsection{Electrostatic Pressure and Mechanical Surface Forces}

\begin{theorem}{Surface Force on a Conducting Boundary}{thm_surf_force}
An infinitesimal surface element \(\dd{A}\) on a charged conductor carrying local surface charge density \(\sigma\) experiences an outward electrostatic force directed along the exterior normal \(\hat{\bm{n}}\):
\begin{equation}
    \dd{\bm{F}} = \frac{\sigma^2}{2\epsilon_0}\,\hat{\bm{n}}\,\dd{A}.
\end{equation}
The resulting outward electrostatic pressure (force per unit area) is:
\begin{equation}
    P_E = \frac{\sigma^2}{2\epsilon_0} = \frac{1}{2}\epsilon_0 E_{\text{surf}}^2.
\end{equation}
\end{theorem}

\begin{proofbox}[Boundary Field Decomposition]
Isolate a small circular patch of area \(\Delta A\) on the conducting boundary. The total field near the surface decomposes into the field generated by the patch \(\bm{E}_{\text{patch}}\) and the field produced by all other charges in the system \(\bm{E}'\):
\begin{equation}
    \bm{E} = \bm{E}_{\text{patch}} + \bm{E}'.
\end{equation}
By Newton's third law, the patch cannot exert a net force on itself; thus, the mechanical force acting on the patch arises solely from \(\bm{E}'\):
\begin{equation}
    \Delta\bm{F} = (\sigma \Delta A)\bm{E}'.
\end{equation}
Modeling the patch locally as an infinite planar sheet with surface charge \(\sigma\):
\begin{equation}
    \bm{E}_{\text{patch}} = \begin{cases}
        +\dfrac{\sigma}{2\epsilon_0}\hat{\bm{n}}, & \text{immediately above (outside)}, \\[1em]
        -\dfrac{\sigma}{2\epsilon_0}\hat{\bm{n}}, & \text{immediately below (inside)}.
    \end{cases}
\end{equation}
The total electric fields immediately above and below the surface are:
\begin{align}
    \bm{E}_{\text{above}} &= \bm{E}' + \frac{\sigma}{2\epsilon_0}\hat{\bm{n}}, \\
    \bm{E}_{\text{below}} &= \bm{E}' - \frac{\sigma}{2\epsilon_0}\hat{\bm{n}}.
\end{align}
Summing these two expressions eliminates \(\bm{E}_{\text{patch}}\):
\begin{equation}
    \bm{E}' = \frac{1}{2}\left(\bm{E}_{\text{above}} + \bm{E}_{\text{below}}\right).
\end{equation}
For an electrostatic conductor in equilibrium, \(\bm{E}_{\text{below}} = \bm{0}\) and \(\bm{E}_{\text{above}} = \frac{\sigma}{\epsilon_0}\hat{\bm{n}}\). Hence:
\begin{equation}
    \bm{E}' = \frac{1}{2}\left(\frac{\sigma}{\epsilon_0}\hat{\bm{n}} + \bm{0}\right) = \frac{\sigma}{2\epsilon_0}\hat{\bm{n}}.
\end{equation}
The net force acting on the patch carrying charge \(\Delta q = \sigma \Delta A\) is therefore:
\begin{equation}
    \Delta\bm{F} = \Delta q \bm{E}' = (\sigma \Delta A)\left(\frac{\sigma}{2\epsilon_0}\hat{\bm{n}}\right) = \frac{\sigma^2}{2\epsilon_0}\hat{\bm{n}}\Delta A.
\end{equation}
\end{proofbox}

\begin{remark}[Equivalence to External Virtual Work]
For a parallel-plate capacitor, the attractive force exerted by one plate on the other is \(F = \frac{\sigma^2 A}{2\epsilon_0} = \frac{\epsilon_0}{2}E^2 A\). The external work required to pull the plates apart from separation $0$ to separation \(d\) at constant charge \(Q\) is:
\begin{equation}
    W_{\text{ext}} = \int_0^d F\,\dd{z} = \left(\frac{1}{2}\epsilon_0 E^2 A\right) d = \left(\frac{1}{2}\epsilon_0 E^2\right) \text{Volume} = U_E.
\end{equation}
This mechanical derivation confirms the field energy density expression \(u_E = \frac{1}{2}\epsilon_0 E^2\).
\end{remark}

\begin{example}[Electrostatic Energy of a Charged Spherical Shell]
Consider a spherical conducting shell of radius \(a\) carrying net charge \(Q\). We compute its electrostatic self-energy through two independent methods:
\begin{enumerate}
    \item \textbf{Volume Integration of Energy Density}:
    The field is \(\bm{E} = \bm{0}\) for \(r < a\) and \(\bm{E}(r) = \frac{Q}{4\pi\epsilon_0 r^2}\hat{\bm{r}}\) for \(r > a\). The energy density in the exterior region is:
    \begin{equation}
        u_E(r) = \frac{1}{2}\epsilon_0 |\bm{E}|^2 = \frac{\epsilon_0}{2}\left(\frac{Q^2}{(4\pi\epsilon_0)^2 r^4}\right) = \frac{Q^2}{32\pi^2\epsilon_0 r^4}.
    \end{equation}
    Integrating over all space using spherical shells of volume element \(\dd{V} = 4\pi r^2 \dd{r}\):
    \begin{equation}
        U = \int_a^\infty u_E(r)(4\pi r^2 \dd{r}) = \frac{Q^2}{8\pi\epsilon_0}\int_a^\infty \frac{\dd{r}}{r^2} = \frac{Q^2}{8\pi\epsilon_0 a}.
    \end{equation}
    \item \textbf{Assembly Work}:
    At an intermediate stage where the shell holds charge \(q\), its surface potential is \(\phi(q) = \frac{q}{4\pi\epsilon_0 a}\). Bringing an extra charge \(\dd{q}\) from infinity requires work \(\dd{W} = \phi(q)\dd{q}\). Thus:
    \begin{equation}
        W = \int_0^Q \frac{q}{4\pi\epsilon_0 a}\,\dd{q} = \frac{Q^2}{8\pi\epsilon_0 a} = \frac{1}{2} Q V,
    \end{equation}
    matching the field energy integral identically.
\end{enumerate}
\end{example}

\subsection{Capacitor Networks and Circuit Combinations}

\subsubsection{Parallel Connection}
In a parallel topology, the corresponding plates of all \(N\) capacitors are directly connected to common electrical nodes. Consequently, the potential difference across every capacitor is identical:
\begin{equation}
    V_1 = V_2 = \dots = V_N = V.
\end{equation}
The total charge \(Q\) drawn from the voltage source is the sum of the individual charges:
\begin{equation}
    Q = \sum_{i=1}^N Q_i = \sum_{i=1}^N C_i V = \left(\sum_{i=1}^N C_i\right) V.
\end{equation}
Setting \(Q = C_{\text{eq}} V\) yields the equivalent parallel capacitance:
\begin{equation}
    C_{\text{eq}}^{\text{parallel}} = \sum_{i=1}^N C_i.
\end{equation}

\subsubsection{Series Connection}
In a series topology, capacitors are cascaded end-to-end such that intermediate nodes are electrically isolated. By charge conservation on these floating segments, each capacitor acquires the same magnitude of charge:
\begin{equation}
    Q_1 = Q_2 = \dots = Q_N = Q.
\end{equation}
The total potential difference across the network is the sum of the individual potential drops:
\begin{equation}
    V = \sum_{i=1}^N V_i = \sum_{i=1}^N \frac{Q}{C_i} = Q \sum_{i=1}^N \frac{1}{C_i}.
\end{equation}
Equating this to \(V = Q / C_{\text{eq}}\) establishes:
\begin{equation}
    \frac{1}{C_{\text{eq}}^{\text{series}}} = \sum_{i=1}^N \frac{1}{C_i}.
\end{equation}

\begin{example}[Composite Network Reduction]
Consider a sub-network where two capacitors \(C_1\) and \(C_2\) are placed in parallel, and this combination is connected in series with a third capacitor \(C_3\).
\begin{enumerate}
    \item First, reduce the parallel pair:
    \begin{equation}
        C_{1,2} = C_1 + C_2.
    \end{equation}
    \item Next, combine \(C_{1,2}\) in series with \(C_3\):
    \begin{equation}
        \frac{1}{C_{\text{eq}}} = \frac{1}{C_{1,2}} + \frac{1}{C_3} = \frac{1}{C_1 + C_2} + \frac{1}{C_3} = \frac{(C_1 + C_2) + C_3}{(C_1 + C_2)C_3}.
    \end{equation}
    Inverting gives the total equivalent capacitance:
    \begin{equation}
        C_{\text{eq}} = \frac{(C_1 + C_2)C_3}{C_1 + C_2 + C_3}.
    \end{equation}
\end{enumerate}
\end{example}

\subsection{Multipole Expansion of Localized Charge Distributions}

\subsubsection{Series Derivation via Legendre Generating Functions}

Consider a bounded, localized continuous charge distribution with density \(\rho(\gr')\) contained entirely within a sphere of radius \(R_{\text{max}}\) centered at the origin. We seek an asymptotic expansion for the electrostatic potential \(\phi(\gr)\) at an observation point \(\gr\) situated in the far-field regime (\(r \gg r'\)).

\begin{center}
\begin{tikzpicture}[scale=1.0, >=Stealth]
    \draw[->, thick] (0,0) -- (4,0) node[right] {\(y\)};
    \draw[->, thick] (0,0) -- (0,3.5) node[above] {\(z\)};
    \draw[->, thick] (0,0) -- (-2,-1.5) node[below left] {\(x\)};
    \node[below right] at (0,0) {\(O\)};
    \draw[rotate around={25:(1,1)}, fill=teal!15, draw=teal!80!black, thick] (1,1) ellipse (1.2 and 0.8);
    \node[teal!80!black] at (0.6, 1.6) {\(\rho(\gr')\)};
    \coordinate (dV) at (1.2, 1.1);
    \draw[fill=blue!60] (dV) rectangle ++(0.2, 0.2);
    \node[above left] at (dV) {\(\dd{V'}\)};
    \draw[->, thick, ForestGreen] (0,0) -- (dV) node[midway, below right] {\(\gr'\)};
    \coordinate (P) at (4.5, 2.8);
    \filldraw (P) circle (1.5pt) node[right] {\(P(\gr)\)};
    \draw[->, thick, RoyalBlue] (0,0) -- (P) node[midway, above left] {\(\gr\)};
    \draw[->, thick, purple] (1.3, 1.2) -- (P) node[midway, above] {\(\gr^* = \gr - \gr'\)};
    \draw[thick] (0.5, 0.31) arc (32:55:0.6);
    \node at (0.7, 0.6) {\(\alpha\)};
\end{tikzpicture}
\end{center}

The exact potential at \(P(\gr)\) is:
\begin{equation}
    \phi(\gr) = \frac{1}{4\pi\epsilon_0} \int_V \frac{\rho(\gr')}{|\gr - \gr'|}\,\dd{V'}.
\end{equation}
Let \(\alpha\) denote the interior angle between the source coordinate \(\gr'\) and the field position \(\gr\), such that \(\gr \cdot \gr' = r r'\cos\alpha\). By the law of cosines:
\begin{equation}
    (r^*)^2 = |\gr - \gr'|^2 = r^2 + (r')^2 - 2rr'\cos\alpha = r^2 \left[ 1 + \left(\frac{r'}{r}\right)^2 - 2\left(\frac{r'}{r}\right)\cos\alpha \right].
\end{equation}
Defining the dimensionless small parameter \(\epsilon \equiv \left(\frac{r'}{r}\right)\left(\frac{r'}{r} - 2\cos\alpha\right) \ll 1\):
\begin{equation}
    \frac{1}{r^*} = \frac{1}{r}(1 + \epsilon)^{-1/2} = \frac{1}{r}\left(1 - \frac{1}{2}\epsilon + \frac{3}{8}\epsilon^2 - \frac{5}{16}\epsilon^3 + \dots\right).
\end{equation}
Regrouping terms in powers of the ratio \((r'/r)\):
\begin{equation}
    \frac{1}{|\gr - \gr'|} = \frac{1}{r}\left[ 1 + \left(\frac{r'}{r}\right)\cos\alpha + \left(\frac{r'}{r}\right)^2 \frac{3\cos^2\alpha - 1}{2} + \left(\frac{r'}{r}\right)^3 \frac{5\cos^3\alpha - 3\cos\alpha}{2} + \dots \right].
\end{equation}
The coefficients of \((r'/r)^n\) correspond to the \textbf{Legendre polynomials} \(P_n(\cos\alpha)\), defined by the generating function:
\begin{equation}
    \frac{1}{\sqrt{1 - 2xt + t^2}} = \sum_{n=0}^\infty t^n P_n(x), \quad \text{where } t = \frac{r'}{r},\; x = \cos\alpha.
\end{equation}

\begin{theorem}{General Multipole Expansion of the Electrostatic Potential}{thm_multipole}
For any localized charge distribution observed at \(r > \max(r')\), the potential expands as a power series in \(1/r\):
\begin{equation}
    \phi(\gr) = \frac{1}{4\pi\epsilon_0} \sum_{n=0}^\infty \frac{1}{r^{n+1}} \int_V (r')^n P_n(\cos\alpha)\rho(\gr')\,\dd{V'}.
\end{equation}
Explicitly writing the first three multipole terms:
\begin{equation}
    \phi(\gr) = \frac{1}{4\pi\epsilon_0}\left[ \frac{1}{r}\int \rho\,\dd{V'} + \frac{1}{r^2}\int r'\cos\alpha\,\rho\,\dd{V'} + \frac{1}{r^3}\int (r')^2 \left(\frac{3\cos^2\alpha - 1}{2}\right)\rho\,\dd{V'} + \dots \right].
\end{equation}
\end{theorem}

\subsubsection{The Monopole and Dipole Moments}

\begin{definition}{Monopole Moment}{monopole_moment}
The \(n = 0\) term represents the \textbf{monopole potential}:
\begin{equation}
    \phi_{\text{mono}}(\gr) = \frac{1}{4\pi\epsilon_0} \frac{Q}{r}, \quad \text{where } Q \equiv \int_V \rho(\gr')\,\dd{V'}.
\end{equation}
The monopole moment is the total net charge \(Q\) (a scalar invariant). At leading order, any charged system with \(Q \neq 0\) resembles a point charge when viewed from a large distance.
\end{definition}

\begin{definition}{Electric Dipole Moment}{dipole_moment}
Using the relation \(r'\cos\alpha = \hat{\bm{r}}\cdot\gr'\), the \(n = 1\) dipole contribution simplifies to:
\begin{equation}
    \phi_{\text{dip}}(\gr) = \frac{1}{4\pi\epsilon_0 r^2} \int_V (\hat{\bm{r}}\cdot\gr')\rho(\gr')\,\dd{V'} = \frac{1}{4\pi\epsilon_0}\frac{\hat{\bm{r}}\cdot\bm{p}}{r^2},
\end{equation}
where the \textbf{electric dipole moment} \(\bm{p}\) is the vector integral:
\begin{equation}
    \bm{p} \equiv \int_V \gr'\rho(\gr')\,\dd{V'}.
\end{equation}
For discrete distributions of point charges \(q_i\) positioned at vectors \(\gr_i'\), the dipole moment is:
\begin{equation}
    \bm{p} = \sum_{i=1}^N q_i \gr_i'.
\end{equation}
\end{definition}

\begin{example}[Physical Dipole along the \(z\)-Axis]
Consider two point charges \(+q\) and \(-q\) positioned on the \(z\)-axis at \(\gr_+ = +\frac{d}{2}\hat{\bm{z}}\) and \(\gr_- = -\frac{d}{2}\hat{\bm{z}}\), respectively. The net dipole moment is:
\begin{equation}
    \bm{p} = q\left(+\frac{d}{2}\hat{\bm{z}}\right) + (-q)\left(-\frac{d}{2}\hat{\bm{z}}\right) = q d\,\hat{\bm{z}}.
\end{equation}
In spherical coordinates (\(\hat{\bm{r}}\cdot\hat{\bm{z}} = \cos\theta\)), the far-field dipole potential evaluates to:
\begin{equation}
    \phi_{\text{dip}}(r, \theta) = \frac{p\cos\theta}{4\pi\epsilon_0 r^2}.
\end{equation}
\end{example}

\subsubsection{The Electric Field of a Pure Dipole}

To evaluate the far-field electric field \(\bm{E}_{\text{dip}} = -\grad \phi_{\text{dip}}\) in spherical coordinates with azimuthal symmetry (\(\pdv*{\phi}{\varphi} = 0\)):
\begin{align}
    E_r &= -\pdv{\phi}{r} = -\pdv{}{r}\left(\frac{p\cos\theta}{4\pi\epsilon_0 r^2}\right) = \frac{2p\cos\theta}{4\pi\epsilon_0 r^3}, \\
    E_\theta &= -\frac{1}{r}\pdv{\phi}{\theta} = -\frac{1}{r}\left(-\frac{p\sin\theta}{4\pi\epsilon_0 r^2}\right) = \frac{p\sin\theta}{4\pi\epsilon_0 r^3}.
\end{align}
Combining both components yields:
\begin{equation}
    \bm{E}_{\text{dip}}(r, \theta) = \frac{1}{4\pi\epsilon_0 r^3}\left(2p\cos\theta\,\hat{\bm{r}} + p\sin\theta\,\hat{\bm{\theta}}\right).
\end{equation}

\begin{proposition}{Coordinate-Free Formulation of the Dipole Field}{coord_free_dipole}
Using the vector identity \(\bm{p} = (\bm{p}\cdot\hat{\bm{r}})\hat{\bm{r}} + (\bm{p}\cdot\hat{\bm{\theta}})\hat{\bm{\theta}} = (p\cos\theta)\hat{\bm{r}} - (p\sin\theta)\hat{\bm{\theta}}\), the dipole electric field can be written in coordinate-free form:
\begin{equation}
    \bm{E}_{\text{dip}}(\gr) = \frac{1}{4\pi\epsilon_0 r^3}\left[ 3(\bm{p}\cdot\hat{\bm{r}})\hat{\bm{r}} - \bm{p} \right].
\end{equation}
\end{proposition}

\subsubsection{Coordinate Origin Invariance and Transformation Laws}

\begin{proposition}{Origin Dependence of the Dipole Moment}{prop_origin_dependence}
Let a coordinate system \(O\) be translated to a new origin \(\widetilde{O}\) displaced by vector \(\bm{a}\). Under this translation, source position vectors transform as:
\begin{equation}
    \widetilde{\gr}' = \gr' - \bm{a}.
\end{equation}
The dipole moment \(\widetilde{\bm{p}}\) evaluated with respect to the new origin is:
\begin{equation}
    \widetilde{\bm{p}} = \int_V \widetilde{\gr}'\rho(\gr')\,\dd{V'} = \int_V (\gr' - \bm{a})\rho(\gr')\,\dd{V'} = \bm{p} - Q\bm{a}.
\end{equation}
\end{proposition}

\begin{remark}[Gauge-like Invariance Condition]
\leavevmode
\begin{itemize}
    \item If the net charge is zero (\(Q = 0\)), then \(\widetilde{\bm{p}} = \bm{p}\). The electric dipole moment is \textbf{strictly origin-independent} if and only if the total charge of the distribution vanishes.
    \item If \(Q \neq 0\), the dipole moment is arbitrary and depends entirely on the choice of coordinate origin. In this case, an origin shift can always be chosen to set \(\widetilde{\bm{p}} = \bm{0}\), placing the origin at the ``center of charge'' \(\bm{R}_c \equiv \frac{1}{Q}\int \gr'\rho\,\dd{V'}\).
\end{itemize}
\end{remark}

\begin{example}[Dipole Moment of Asymmetric Point Charges under Origin Shifts]
Consider two point charges \(q_1 = 3q\) and \(q_2 = -q\) separated by a distance \(a\). The total charge is \(Q = 3q - q = 2q \neq 0\).
\begin{enumerate}
    \item \textbf{Origin at \(-q\) with $3q$ at \(z = a\)}:
    \begin{equation}
        \bm{p}_{(a)} = (-q)(\bm{0}) + (3q)(a\hat{\bm{z}}) = 3qa\,\hat{\bm{z}} \implies \phi(r, \theta) \simeq \frac{1}{4\pi\epsilon_0}\left[\frac{2q}{r} + \frac{3qa\cos\theta}{r^2}\right].
    \end{equation}
    \item \textbf{Origin at \(+3q\) with \(-q\) at \(z = -a\)}:
    The origin is displaced along \(\bm{a} = a\hat{\bm{z}}\) relative to case (a):
    \begin{equation}
        \bm{p}_{(b)} = (3q)(\bm{0}) + (-q)(-a\hat{\bm{z}}) = qa\,\hat{\bm{z}}.
    \end{equation}
    Verifying with the transformation law:
    \begin{equation}
        \widetilde{\bm{p}} = \bm{p}_{(a)} - Q\bm{a} = 3qa\hat{\bm{z}} - (2q)(a\hat{\bm{z}}) = qa\,\hat{\bm{z}}.
    \end{equation}
    The resulting far-field potential in this frame is:
    \begin{equation}
        \phi(r, \theta) \simeq \frac{1}{4\pi\epsilon_0}\left[\frac{2q}{r} + \frac{qa\cos\theta}{r^2}\right].
    \end{equation}
    \item \textbf{Origin at \(-q\) with $3q$ at \(y = a\)}:
    \begin{equation}
        \bm{p}_{(c)} = (-q)(\bm{0}) + (3q)(a\hat{\bm{y}}) = 3qa\,\hat{\bm{y}}.
    \end{equation}
    Since \(\hat{\bm{y}}\cdot\hat{\bm{r}} = \sin\theta\sin\varphi\), the expansion becomes:
    \begin{equation}
        \phi(r, \theta, \varphi) \simeq \frac{1}{4\pi\epsilon_0}\left[\frac{2q}{r} + \frac{3qa\sin\theta\sin\varphi}{r^2}\right].
    \end{equation}
\end{enumerate}
\end{example}

\subsubsection{Higher-Order Moments: The Quadrupole Potential}

When both the total charge \(Q\) and the net dipole moment \(\bm{p}\) vanish, the potential is dominated at large distances by the \(n = 2\) \textbf{quadrupole term}:
\begin{equation}
    \phi_{\text{quad}}(\gr) = \frac{1}{4\pi\epsilon_0 r^3} K_2(\hat{\bm{r}}), \quad \text{where } K_2(\hat{\bm{r}}) \equiv \int_V (r')^2 \left(\frac{3\cos^2\alpha - 1}{2}\right)\rho(\gr')\,\dd{V'}.
\end{equation}

\begin{example}[Non-Uniform Sphere: Vanishing Monopole and Dipole Moments]
Consider a sphere of radius \(R\) centered at the origin with charge density \(\rho(r',\theta') = k\frac{R}{(r')^2}(R - 2r')\sin\theta'\)[cite: 22]. For observation points along the \(z\)-axis (\(\gr = z\hat{\bm{z}}\), \(\dd{V'} = (r')^2\sin\theta'\dd{r'}\dd{\theta'}\dd{\varphi'}\)):
\begin{itemize}
    \item \textbf{Monopole Moment (\(Q\))}: The radial integral vanishes identically[cite: 22]:
    \begin{equation}
        Q = kR \int_0^{2\pi}\dd{\varphi'} \int_0^\pi \sin^2\theta'\dd{\theta'} \int_0^R (R - 2r')\,\dd{r'} = 0 \implies \phi_{\text{mono}} = 0.
    \end{equation}
    \item \textbf{Dipole Moment (\(\bm{p}\))}: Azimuthal symmetry ensures \(p_x = p_y = 0\), while the angular \(\theta'\)-integral vanishes for the \(z\)-component[cite: 23]:
    \begin{equation}
        p_z = \int z'\rho\,\dd{V'} \propto \int_0^\pi \sin^2\theta'\cos\theta'\,\dd{\theta'} = \left[\frac{\sin^3\theta'}{3}\right]_0^\pi = 0 \implies \bm{p} = \bm{0}.
    \end{equation}
\end{itemize}
\end{example}

\begin{example}[Non-Uniform Sphere: Leading Quadrupole Term]
Because both \(Q = 0\) and \(\bm{p} = \bm{0}\), the lowest non-vanishing contribution along the \(z\)-axis (\(\alpha = \theta'\)) is the quadrupole term \(K_2\)[cite: 22, 23]:
\begin{equation}
    K_2 = \frac{kR}{2} \underbrace{\int_0^{2\pi}\dd{\varphi'}}_{2\pi} \underbrace{\int_0^\pi (3\cos^2\theta' - 1)\sin^2\theta'\,\dd{\theta'}}_{-\pi/8} \underbrace{\int_0^R (Rr'^2 - 2r'^3)\,\dd{r'}}_{-R^4/6} = \frac{k\pi^2 R^5}{48}.
\end{equation}
The far-field potential for \(z \gg R\) is therefore dominated by the quadrupole decay[cite: 23]:
\begin{equation}
    \phi(z) \simeq \frac{1}{4\pi\epsilon_0}\frac{K_2}{z^3} = \frac{1}{4\pi\epsilon_0}\frac{k\pi^2 R^5}{48 z^3}.
\end{equation}
\end{example}

\subsection{Summary: Foundational Framework of Electrostatic Methods}

\begin{center}
\renewcommand{\arraystretch}{1.4}
\begin{tabular}{|p{4.2cm}|p{6.0cm}|p{4.8cm}|}
\hline
\textbf{Method} & \textbf{Governing Equations} & \textbf{Typical Applications} \\ \hline
1. \textbf{Coulomb Superposition} & \(\displaystyle \bm{E}(\gr) = \frac{1}{4\pi\epsilon_0}\int \frac{\rho(\gr')}{|\gr - \gr'|^2}\hat{\bm{r}}^*\,\dd{V'}\) & Localized discrete or continuous charges \\ \hline
2. \textbf{Gauss's Flux Law} & \(\displaystyle \oint_{\partial V} \bm{E} \cdot \dd{\bm{a}} = \frac{Q_{\text{enc}}}{\epsilon_0}\) & High symmetry (spherical, cylindrical, planar) \\ \hline
3. \textbf{Electric Potential} & \(\displaystyle \bm{E} = -\grad\phi, \quad \phi(\gr) = \frac{1}{4\pi\epsilon_0}\int \frac{\rho(\gr')}{|\gr - \gr'|}\,\dd{V'}\) & Direct scalar integration over sources \\ \hline
4. \textbf{Poisson \& Laplace} & \(\displaystyle \nabla^2\phi = -\frac{\rho}{\epsilon_0}, \quad \nabla^2\phi = 0\) & Boundary-value problems with conductors \\ \hline
5. \textbf{Uniqueness Theorems} & \(\displaystyle \nabla^2\phi = -\frac{\rho}{\epsilon_0} \text{ with } \phi|_{\partial V} \text{ or } Q_i \text{ fixed}\) & Rigorous mathematical uniqueness of fields \\ \hline
6. \textbf{Method of Images} & Replace conductors with virtual charges & Point charges near planes or spheres \\ \hline
7. \textbf{Multipole Expansion} & \(\displaystyle \phi(\gr) = \frac{1}{4\pi\epsilon_0}\left[\frac{Q}{r} + \frac{\bm{p}\cdot\hat{\bm{r}}}{r^2} + \frac{K_2}{r^3} + \dots\right]\) & Far-field approximations (\(r \gg \text{size}\)) \\ \hline
\end{tabular}
\end{center}